\section{Local changes of scale}
\label{sec:renormalization}

All averaging in this section is in the critical nearest-neighbor reference
measure.  We change the interaction functions on its original spins,
preserving the partition integral exactly at each step.  The linear part
is conditional averaging.  The nonlinear corrections form a connected-list
expansion, whose bounds are uniform in the volume.  The
radial estimates will then show that the linear map contracts homogeneous
backgrounds after one thermal direction is removed.  None of these
estimates presupposes correlation bounds in an interacting model.

Fix a large integer $L$ and write
\[
 r_j=r_0L^j,\qquad r=r_j,\qquad R=Lr,\qquad
 M=\lceil(\log L)^2\rceil.
\]
The order of choices is important: $L$ is chosen first, $r_0$ is then made
large enough for the fixed-ratio estimates of the preceding section, and
the permissible perturbation is finally made small.  Constants denoted by
$C_L$ may depend on the fixed value of $L$ but not on $j$ or the volume.

\subsection{Carriers, constants, and norms}

An interaction is a list $h=(f_{A,z})$.  Its carrier $A$ is a finite connected
set of lattice sites, its anchor $z$ belongs to $A$, and $f_{A,z}$ depends
only on the spins available in $A$.  In a domain a carrier may extend
outside the domain; those exterior sites are not additional variables.
On a torus we use connected torus carriers and the torus distance.  Terms
at an identical index $(A,z)$ are added as functions.  We do not require a
unique decomposition of a function into different carriers.

For backgrounds, functions are considered modulo constants.  We choose
their canonical representatives by subtracting their expectation under
independent uniform values of the available variable spins.  Denote this
projection by $\pi$.  It is linear, preserves the relevant symmetries, and
has sup-norm operator bound at most two.  The removed scalar is retained
whenever sources are present.  Thus, in a finite volume, an identity
\[
 \E_0 e^{H}=e^c\E_0 e^{H'}
\]
includes the accumulated scalar $c$; source-dependent constants are never
allowed to reenter $H'$ in a later transformation.  Background constants
cancel upon partition normalization.

For a carrier $A$, let $A+B_r$ be its max-norm $r$-neighborhood and put
\begin{equation}
 N_r(A)=\frac{|A+B_r|}{r^2},\qquad
 \|h\|_r=\sup_x
 \sum_{\substack{A\cap B_r(x)\ne\varnothing\\z\in A}}
       \|f_{A,z}\|_\infty e^{N_r(A)}.
 \label{rg:norm}
\end{equation}
Lattice counting measure is used in $N_r$.  The harmless bounded changes
caused by rounding radii are absorbed in constants.  These componentwise
summation norms give Banach spaces, and the translation, square-symmetry,
and spin-flip restrictions define closed subspaces.  The homogeneous
background space is translation invariant, invariant under the square
lattice symmetries, and spin-flip even.

A source coefficient is supplied with a nonempty set of labels, and its
carrier must contain their marked sites.  For any one fixed label set its
norm is the total-list norm
\begin{equation}
 \|u\|_{r,\bullet}=
 \sum_{A,z}\|u_{A,z}\|_\infty e^{N_r(A)}.
 \label{rg:marked-norm}
\end{equation}
The anchor of a singleton-source list is its marked site.  For other
marked lists one may choose a mark according to a fixed rule.  The
background norm has a volume factor when all anchors are summed in a
block; the marked norm does not.

We call a carrier short if $N_r(A)\le M$.  A connected carrier satisfies
$\diam A\le CrN_r(A)$: along a path between distant points, select
separated $r$-boxes, each contributing order $r^2$ to the thickened
carrier.  Consequently every short carrier anchored at $z$ is contained
in $B_{CMr}(z)\subset B_R(z)$ once $L$ is large.  All carriers of the
distinguished fields below are short.

\subsection{Geometry and the exact analytic map}

Give an input carrier the protective neighborhood
\[
 U(A)=A+B_{2R}.
\]
Two occurrences in a list are joined when their protective neighborhoods
intersect.  Connectedness always refers to this overlap graph when it is
applied to a list of interactions.

\begin{lemma}[Carrier geometry]
\label{rg:geometry}
For connected $A$ and $R=Lr$, with $L$ sufficiently large,
\begin{equation}
 N_R(U(A))\le C\bigl(1+L^{-1}N_r(A)\bigr).
 \label{rg:geometry-eq}
\end{equation}
Moreover $A+B_{5R}$ can be covered by at most
$C_L(1+N_r(A))$ squares of radius $r$.  These estimates hold on a torus
while $R$ is less than a fixed sufficiently small fraction of its side.
\end{lemma}

\begin{proof}
Choose a maximal collection of mutually separated points of $A$, with
separation a fixed multiple of $R$.  Enlarged $R$-boxes about them cover
$U(A)+B_R$, so $N_R(U(A))$ is bounded by a constant times their number.
If there is only one such point this gives the constant term in
\eqref{rg:geometry-eq}.  Otherwise, from each chosen point a path in $A$
must leave a smaller concentric box.  The $r$-neighborhood of the initial
part of that path contains at least $cRr$ sites.  Choose these smaller
boxes disjoint, including the $r$-thickening.  Summing their contributions
to $|A+B_r|$ proves the other term.  The argument uses only boxes of size
comparable to $R$ and applies unchanged in a torus chart.  A coarse
$r$-grid covering of $A+B_{5R}$ costs at most $C_L$ times a covering of
$A+B_r$, proving the second assertion.
\end{proof}

The linear rule sends a short input function to $P_{R,z}f_{A,z}$ on
carrier $B_R(z)$, and leaves a long input unchanged.  At this point
constants have not yet been removed.  Here and below the cutting contour
is included in the square, as in \Cref{model:reference}.

Start with finitely many input terms.  The connected coefficient
estimates below will define the same local rule for arbitrary lists
in the rooted-norm ball.  To define the nonlinear rule, expand the
exponential of the input interaction
in ordered lists, with coefficient $1/k!$ for a list of
length $k$.  On each singleton component of its overlap graph apply the
linear rule.  On every component with at least two vertices retain the
product of its original input functions.  Multiply the component
weights and sum the resulting series, denoted by $\mathcal W_j(h)$.
This is a function of the reference spins, with constant term one in
$h$.  Take its logarithm pointwise, before integrating
$\mathcal W_j(h)$ against the reference measure.

The coefficient involving two input occurrences illustrates the rule.
Write $g_i=P_{R,z_i}f_i$ for a short occurrence and $g_i=f_i$ for a
long one.  With temporary parameters $s,t$, the coefficient of $st$ in
$\mathcal W_j(sf_1+tf_2)$ is $g_1g_2$ when the protective neighborhoods
are disjoint and $f_1f_2$ when they intersect.  Therefore, before removing
constants,
\begin{equation}\label{rg:pair-coefficient}
 [st]\log\mathcal W_j(sf_1+tf_2)
 =\begin{cases}
   0,&U(A_1)\cap U(A_2)=\varnothing,\\
   f_1f_2-g_1g_2,&U(A_1)\cap U(A_2)\ne\varnothing.
  \end{cases}
\end{equation}
Here $[st]$ means extraction of that power-series coefficient.  Thus
separated occurrences acquire no new interaction, while the correction
for an overlapping pair retains its original product.  The proof below
establishes this connectedness property at every order.

At order at least two assign the connected coefficient to the carrier
\begin{equation}
 Q=\bigcup_{i=1}^k U(A_i).
 \label{rg:test-carrier}
\end{equation}
Average the choice of its output anchor over its input anchors.  This
preserves translation and square symmetries.  For marked coefficients
use the chosen marked-anchor convention instead.  Unions may be enlarged
by a bounded number of lattice sites when needed to make the carrier
connected.  Finally apply $\pi$ and record the removed scalars.  Write
$\mathcal R_j$ for the resulting nonscalar map and $T_j$ for its derivative
at zero.

For estimates it is useful to retain $Q$ as an auxiliary \emph{test
carrier} even at order one: there $Q=U(A)$, although the actual output
carrier is $B_R(z)$ or $A$.  This auxiliary set is used only to estimate
locality; it does not enlarge the actual carrier of a linearly averaged
term.

For source bookkeeping, attach a distinct formal label to each marked
insertion and impose $t_\ell^2=0$ for each label $\ell$. If $t_S$
denotes the product of the labels in $S$, set
$t_St_T=t_{S\cup T}$ when $S\cap T=\varnothing$, and $t_St_T=0$
otherwise. Thus every retained monomial uses each label at most once.
The full finite source algebra is specified in \Cref{sec:correlations}.

\begin{proposition}[Exact analytic transformation]
\label{rg:map}
For each fixed sufficiently large $L$, the maps $\mathcal R_j$ are analytic
on a common complex neighborhood of zero in \eqref{rg:norm}, uniformly in
$j$ and the finite volume.  They satisfy $\mathcal R_j(0)=0$ and
\begin{equation}
 \|\mathcal R_j(h)-T_jh\|_R\le C_L\|h\|_r^2,\qquad
 \|D\mathcal R_j(h)-T_j\|\le C_L\|h\|_r.
 \label{rg:quadratic}
\end{equation}
The derivative bound holds also from the marked norm at $r$ to the marked
norm at $R$.  For a marked coefficient, the corresponding scalar-increment
functional is bounded, and its derivative differs from its derivative
at zero by at most
$C_L\|h\|_r$ in the same marked norms.

The coefficient estimates remain valid if the output weight
$e^{N_R(\text{output carrier})}$ is replaced by $e^{2N_R(Q)}$.  In
particular, restricting an output list to connected test carriers
with diameter at least $D$ gives the additional factor
\begin{equation}
 e^{-cD/R}.
 \label{rg:diameter-gain}
\end{equation}
This factor can be retained simultaneously with the norm of any small
input, including a marked remainder or a background.  The assertion holds
before scalar projection as well as after it, and includes long unchanged
singletons with their test neighborhoods.

For any fixed number of square-free source labels, the induced maps on
the corresponding coefficients are polynomial and uniformly Lipschitz
on bounded sets of marked coefficients, while the background remains
in this common small ball.  All the preceding enhanced-weight and
diameter bounds apply to these maps.  Constants may depend on the number
of labels and their norm bounds.

In finite volume the transformation, including its scalar factor,
preserves the partition integral exactly.  The same conclusions apply to
a preliminary transformation of finite-range microscopic interactions
to a fixed large scale $r_0$, averaging every singleton there; its
constants may also depend on $r_0$ and the finite microscopic range.
Finitely many prescribed source carriers may instead be left unchanged
if they are too large for these cutting squares, with constants allowed
to depend on those carriers.
\end{proposition}

\begin{proof}
We prove the integral identity first, then bound the connected
coefficients and sum them in the weighted norms.  The same estimates
will also control marked coefficients and the removed scalars.

\emph{Integral preservation.}
If an occurrence is an isolated vertex of a list's overlap graph, its
cutting square has a nearest-neighbor buffer from the supports of all
other factors.  Cutting squares of distinct isolated vertices likewise
have such a buffer from one another.  Conditional on their cutting
contours and the remaining spins, the reference measure makes their
interiors independent.  Replacing their factors by conditional
expectations therefore leaves the list's integral unchanged.  A long
singleton is not altered.  This proves exactness for every list, including
lists with nonsingleton components; it does not require the components
themselves to be independent in the unconditioned reference measure.
Consequently the integral of $\mathcal W_j(h)$ is the original partition
integral.  Exponentiating the scalar and nonscalar parts of
$\log\mathcal W_j(h)$ recovers $\mathcal W_j(h)$ pointwise, giving
the required partition identity.

\emph{Connected coefficients.}
Fix $k$ labeled occurrences and
their overlap graph $G$.  The transformed moment of any subset of these
labels factors over the components of its induced graph.  It follows
either by the partition formula for a logarithm, or by the following
polymer expansion, that the logarithmic coefficient vanishes unless
$G$ is connected.  For connected $G$ its sup norm is at most
\begin{equation}
 2^{k-1}\tau(G)\prod_{i=1}^k\|f_i\|_\infty,
 \label{rg:tree-coefficient}
\end{equation}
where $\tau(G)$ is the number of spanning trees of $G$.

The connected logarithm uses the hard-core Mayer expansion and the
Penrose tree-graph bound \cite[Section~2 and Proposition~5]{FP2007}.
We use the minimum-spanning-tree partition scheme
\cite[Section~3.2]{PY2017}; the extra factor $2^{k-1}$ will come
from the internal/external edge markings below.
To prove \eqref{rg:tree-coefficient}, use square-free variables for the
$k$ occurrences.  Regard a connected subset of labels as a polymer.  Its
activity is the product of the original functions if its size is at
least two, and the linearly transformed function if its size is one.
Polymers are incompatible when their label sets overlap or an edge of
$G$ joins them.  The hard-core polymer partition polynomial is exactly
the transformed moment polynomial, because every induced graph has a
unique decomposition into its components.

Its logarithm is a sum over partitions of the labels into connected
blocks.  For a fixed partition, the absolute Mayer coefficient is bounded
by the spanning-tree sum of the quotient contact graph.  For completeness,
order that graph's edges strictly and group its connected spanning
subgraphs according to their greedy minimum spanning tree.  For a given
tree, the possible extra edges are optional independently: an edge may be
added precisely when it is later than every edge on its tree path.  Their
alternating sum is zero or has absolute value one.  Summing over trees
proves the bound.

Choose a spanning tree within each connected block, and for every edge
of a quotient tree choose an original contact edge.  Allowing all these
choices only increases the bound.  Their union is a spanning tree of
$G$ whose edges are marked as internal to a block or between blocks.
Conversely, the tree and its edge marks determine the blocks and these
choices.  There are at most $2^{k-1}$ markings for each spanning tree of
$G$.  Finally each polymer activity is bounded by the product of its
input norms, since conditional averaging is a sup-norm contraction.
This proves \eqref{rg:tree-coefficient}.  Repeated input indices cause no
change: their occurrences were labeled, and the original expansion still
carries $1/k!$.

\emph{Weighted summation.}
Subadditivity and \Cref{rg:geometry} give
\begin{equation}
 2N_R(Q)\le 2\sum_iN_R(U(A_i))
       \le Ck+\frac{C}{L}\sum_iN_r(A_i).
 \label{rg:weight-allocation}
\end{equation}
Choose $L$ so large that the last coefficient is, for example, less than
one quarter.  Thus a fixed fraction of each input exponential weight
remains after paying the enhanced output weight.

Root a spanning tree at an input whose test neighborhood meets the
output rooted block.  Such a root carrier meets $B_{CR}(x)$; covering that
square by $O(L^2)$ input rooted blocks costs $C_L$.  If a child touches its
parent's protective neighborhood, its carrier meets
$A_{\rm parent}+B_{5R}$.  By \Cref{rg:geometry}, the sum of possible child
terms costs at most $C_L(1+N_r(A_{\rm parent}))$ times the child norm.
For a vertex with $b$ children, the unused exponential weight pays this
factor to the power $b$, using
\[
 (1+t)^b e^{-ct}\le b!\,C_c^{\,b+1},\qquad t\ge0.
\]
Reserve fixed fractions of the weight for this inequality and for
\eqref{rg:weight-allocation}; they do not depend on $k$.

The resulting tree bound is $C_L^k\prod_i b_i!$ times the product of
input norms.  Summing rooted labeled trees and dividing by $k!$ costs
only another constant to the power $k$: ordering the children identifies
these objects with plane rooted trees, whose number on $k$ vertices is
the Catalan number $\frac1k\binom{2k-2}{k-1}\le4^{k-1}$.
Choices of root and the factor $2^{k-1}$ are included in $C_L^k$.
Consequently the order-$k$ contribution, including the doubled test
weight, is bounded by
\begin{equation}
 C_L^k\|h\|_r^k.
 \label{rg:order-bound}
\end{equation}
The order-one rule is also bounded with this weight.  For its short part
the output test neighborhoods have bounded $N_R$ and the root covering
cost is $C_L$; for its long part use
\eqref{rg:weight-allocation} directly.  Projection costs at most two.
Summing \eqref{rg:order-bound} proves analyticity and
\eqref{rg:quadratic}, including the usual Lipschitz bounds on a smaller
ball.

If a connected test carrier has diameter at least $D$, its
$R$-neighborhood has cardinality at least $cDR$, hence
$N_R(Q)\ge cD/R$.  The estimate with $e^{2N_R(Q)}$, after retaining the
ordinary output weight, leaves \eqref{rg:diameter-gain}.  In the torus the
same argument applies to distances up to a fixed fraction of the side;
a wrapping or non-chart connected carrier in particular crosses such a
distance.  For a singleton the same estimate is made on $U(A)$ even when
its actual carrier is smaller.  Each input norm has remained a separate
factor throughout the argument.  A small factor in one of them therefore
persists when the diameter restriction is imposed.

\emph{Marked coefficients and scalars.}
For marked coefficients, root a tree at an occurrence carrying a
specified mark.  Summation of this root is by
\eqref{rg:marked-norm}, and the remaining argument is identical.
For multiple labels, regard each input coefficient as carrying the
product of its labels and retain only disjoint label sets.  There are
only finitely many allocations for any fixed total label set.  The
background occurrences are summed by the convergent series above;
the marked occurrences give polynomials in the marked coefficients.
Equivalently one can differentiate on a small source polydisc whose
radius depends on their norm bound.  This proves all the marked analytic
and Lipschitz statements.  In particular, terms in a one-mark derivative
other than its zero-background linear term contain a background factor,
which proves the second bound in \eqref{rg:quadratic} in marked norm.
The removed scalar of a function has absolute value at most its sup
norm.  Before discarding its carrier, apply exactly the same estimate;
this proves the asserted bounds for scalar increments and their
background differences.  No parity assumption was needed for these
marked estimates.

For finitely many input terms the partition identity follows first as an
identity of power series near zero.  The convergent logarithmic series
then defines its exponential throughout the common small ball, and
analytic continuation gives the identity there.  In a finite system the
total absolute norm is finite, bounded by finitely many rooted norms:
every nonconstant term meets a variable site.  Truncating the list and
using the proved convergence gives the identity for all interactions in
the ball.  This argument does not expand a thermodynamic partition
function under the interacting measure.

For the preliminary finite-range transformation, take protective
neighborhoods at the fixed radius $r_0$ and make $r_0$ larger than every
microscopic support.  The same exactness and tree proof apply.  All
covering costs are finite constants depending on this radius and the
range; the enhanced output weight costs at most a fixed factor per input
occurrence.  A prescribed larger source carrier is left unchanged in the
singleton rule, with its own protective neighborhood retained.  The
marked rooted estimates then have finite additional constants depending
on that carrier.  This proves the last assertion.
\end{proof}

One useful consequence isolates the long linear terms.  Let
$T_j^{\rm long}$ denote the rule that leaves terms with $N_r(A)>M$
unchanged.  Comparing the input and output weights and covering an
$R$-rooted block by $r$-blocks gives
\begin{equation}
 \|T_j^{\rm long}h\|_R
 \le C L^2 e^{-cM}\|h\|_r=o_L(1)\|h\|_r.
 \label{rg:long-tail}
\end{equation}
The same conclusion holds for a marked list, also after multiplication
by either $L$ or $L^{1/8}$.  The constants $c>0$ and the polynomial
covering factor can be chosen independently of sufficiently large $L$;
the quadratic map constants $C_L$ need not have this property.

\subsection{Distinguished local fields}

Use the microscopic observables
\[
 \mathcal O_1(z)=e(z)=\frac14\sum_{w:\,|w-z|_1=1}\sigma_z\sigma_w,
 \qquad
 \mathcal O_{1/8}(z)=\sigma_z.
\]
For a function in a bulk chart let $m(F)$ denote its plane reference mean,
and set
\[
 d_R(F)=\E^+_{\operatorname{disk}(4R)}F-m(F).
\]
Both quantities here concern the reference model.  Scalar changes to
$F$ do not change $d_R(F)$.

\begin{lemma}[Distinguished fields]
\label{rg:fields}
On carrier $B_{r_j}(z)$ define
\begin{equation}
 v_j^{(d)}(z)=r_j^d\pi P_{r_j,z}\mathcal O_d(z),
 \qquad d\in\{1,1/8\}.
 \label{rg:distinguished}
\end{equation}
For spin the projection removes no constant.  If $r_0$ is sufficiently
large, there are constants $0<c<C<\infty$, independent of $j$ and of the
large blocking ratio $L$, such that
\begin{equation}
 c\le\|v_j^{(d)}(z)\|_\infty\le C.
 \label{rg:field-norms}
\end{equation}
Their single-term marked norms are likewise bounded above and below.
For $R/r_j$ sufficiently large,
\begin{equation}
 c(r_j/R)^d\le |d_R(v_j^{(d)}(z))|
                    \le C(r_j/R)^d,
 \label{rg:field-response}
\end{equation}
where the disk is centered at $z$.
\end{lemma}

\begin{proof}
\Cref{rad:calibration} gives constants $\kappa_d\ne0$ such that
\begin{equation}
 R^d d_R(\mathcal O_d(z))\longrightarrow\kappa_d.
 \label{rg:micro-response}
\end{equation}
Whenever the conditioning square is inside the disk, the Markov property
and the tower property give
\begin{equation}
 d_R(P_{s,z}\mathcal O_d(z))=d_R(\mathcal O_d(z)).
 \label{rg:tower-response}
\end{equation}
This proves \eqref{rg:field-response} for all sufficiently large radii,
and gives the lower bound in \eqref{rg:field-norms} by taking a disk of
radius comparable to $r_j$.  The difference of two expectations of a
canonical representative is bounded by twice its sup norm.

For the upper bound, choose a fixed large ratio $B$, independently of the
later blocking ratio $L$.  At all sufficiently large $s$, apply
\Cref{rad:operator} to
$F=P_{s,z}\mathcal O_d(z)-m(\mathcal O_d)$ and to the outer radius $Bs$.
The relevant remainder exponent is $p=3$ for the square-symmetric even
field and $p=1+1/8$ for the odd field.  Equations
\eqref{rg:micro-response}--\eqref{rg:tower-response} yield
\[
 (Bs)^d\|P_{Bs,z}\mathcal O_d(z)-m(\mathcal O_d)\|_\infty
 \le C+C B^{d-p}
       s^d\|P_{s,z}\mathcal O_d(z)-m(\mathcal O_d)\|_\infty.
\]
Take $B$ large enough that the last coefficient is below one.  Iterating
from radii in a fixed compact interval above the threshold proves a
uniform bound for every large radius.  Integer radii are handled by
ratios tending to $B$, which are allowed in \Cref{rad:operator}.  The
initial compact interval costs only a finite constant.  Passing from
plane centering to $\pi$ costs at most two.  This bound was obtained with
$B$ fixed, so is independent of the later choice of $L$.
Finally $N_{r_j}(B_{r_j}(z))$ is bounded above and below by absolute
constants, giving the marked-norm assertion.
\end{proof}

Let $\boldsymbol v_j$ be the homogeneous thermal list with one term
\begin{equation}
 r_j^{-2}v_j^{(1)}(z)
 \label{rg:thermal-list}
\end{equation}
at every anchor $z$.  Its background norm is comparable to one: a rooted
$r_j$-block meets order $r_j^2$ of these carriers.  Nested conditional
expectations compose, so the linear rule satisfies the exact identities
\begin{equation}
 T_j\boldsymbol v_j=L\boldsymbol v_{j+1},\qquad
 L^dT_jv_j^{(d)}(z)=v_{j+1}^{(d)}(z)
 \quad\text{in the marked space}.
 \label{rg:line-transport}
\end{equation}
All these equalities are modulo constants where appropriate.

\subsection{Linear contraction estimates}

The background and marked norms count local fields differently.  A
homogeneous background has order $L^2$ input blocks per output block,
so its contraction requires a radial remainder with dimension greater
than two.  The right-angle invariant estimate in \Cref{rad:operator}
provides dimension three after the energy response is removed.  A
marked coefficient has no such volume factor: its prescribed
normalization contributes only $L$ for energy or $L^{1/8}$ for spin.
The respective remainder dimensions two and $9/8$ therefore suffice.
At the physical boundary there are order $L$ blocks along the edge,
and the boundary remainder has dimension two.  The following proof
keeps the additional powers of the short-carrier cutoff $M$; since
$M=\lceil(\log L)^2\rceil$, they do not spoil these strict gains.

\begin{proposition}[Bulk, marked, and boundary linear bounds]
\label{rg:linear}
Given $\rho>0$, choose $L$ sufficiently large and then $r_0$ sufficiently
large.  In the homogeneous background space there are uniformly bounded
splittings into the line $\R\boldsymbol v_j$ and a complement such that
\begin{equation}
 T_j=
 \begin{pmatrix}L&B_j\\0&D_j\end{pmatrix},\qquad
 \|D_j\|\le\rho,\qquad \|B_j\|\le C_L.
 \label{rg:background-matrix}
\end{equation}
The coordinates in this splitting are denoted by $(t_j,w_j)$.

For a singleton marked source, in the even sector with $d=1$ or the odd
sector with $d=1/8$, the scaled operator $L^dT_j$ has a corresponding
splitting about $v_j^{(d)}$ with matrix
\begin{equation}
 \begin{pmatrix}1&B_j^{(d)}\\0&D_j^{(d)}\end{pmatrix},\qquad
 \|D_j^{(d)}\|\le\rho,\qquad \|B_j^{(d)}\|\le C_L.
 \label{rg:marked-matrix}
\end{equation}
The even projections can be chosen invariant under square rotations.

In a finite domain consider a boundary family: every carrier touches a
fixed bounded lattice neighborhood of the physical edge.  Its ordinary
linear image, modulo constants, has norm at most $\rho$ times its input
norm, provided the local graph and flatness conditions of
\Cref{rad:boundary} hold on the required windows at the current scale.
For free boundary conditions the family must be even; for either fixed
boundary condition no parity restriction is imposed.  The image is
again a boundary family.
\end{proposition}

\begin{proof}
We first prove the quotient estimate in the homogeneous plane space.
Sum all short input functions at one anchor $z$ into $F_z$.  It is
square-symmetric and even, is supported in $B_s(z)$ with $s=CMr$, and
obeys
\begin{equation}
 \|F_z\|_\infty\le Cr^{-2}\|h\|_r.
 \label{rg:anchor-volume}
\end{equation}
Indeed, every carrier contains its anchor.  For any $x$, all terms
anchored at $z\in B_r(x)$ contribute to the rooted sum in
\eqref{rg:norm}.  Translation invariance therefore makes this sum at
least $|B_r|$ times the total weighted sum at a single anchor.  This
proves \eqref{rg:anchor-volume} without a further factor of $M$.

Subtract $\alpha r^{-2}v_j^{(1)}(z)$ from $F_z$, choosing $\alpha$ to
make its disk response at radius $R$ vanish.  Equations
\eqref{rg:field-response} and \Cref{rad:operator} show
\[
 |\alpha|\le CM\|h\|_r.
\]
The adjusted function remains supported in $B_s(z)$ and has norm at
most $C(1+M)r^{-2}\|h\|_r$.  All its short terms are mapped to the
same carrier $B_R(z)$ and are added there.  Applying
\Cref{rad:operator} with $(d,p)=(1,3)$ to this sum, and then projecting
off constants, bounds the remaining function by
\[
 C(CM/L)^3(1+M)r^{-2}\|h\|_r.
\]
There are $O(R^2)$ anchors whose output squares meet an $R$-rooted block,
and their output weights are bounded constants.  The quotient norm of
the short output is consequently at most
\begin{equation}
 C L^2(1+M)(CM/L)^3\|h\|_r=o_L(1)\|h\|_r.
 \label{rg:bulk-quotient}
\end{equation}
Add the long-term bound \eqref{rg:long-tail}.  Since the distinguished
line is exactly invariant, the same estimate applies after subtracting
any multiple of the input line.  Optimizing gives an operator bound
$o_L(1)$ on the quotient spaces, rather than merely a bound on a
particular representative.

The norm bounds for $\boldsymbol v_j$ give uniformly bounded projections
onto these lines: choose a norm-bounded dual functional equal to one on
$\boldsymbol v_j$.  Choose the functionals first on real lists and
extend them complex-linearly.  Their norms remain uniformly bounded,
and the resulting projections commute with complex conjugation.
The corresponding complement projections have
uniform norm.  In these splittings the lower-left entry vanishes by
\eqref{rg:line-transport}; the quotient estimate bounds the lower-right
entry, and the bounded full linear map bounds the upper-right entry.
Taking $L$ larger if necessary proves
\eqref{rg:background-matrix}.

For a singleton-source list sum its short terms at the marked anchor.
The sum has norm at most $\|u\|_{r,\bullet}$ and support in $B_{CMr}$.
The same response cancellation subtracts at most
$CM^d\|u\|_{r,\bullet}$ times $v_j^{(d)}$.  There is now no volume
factor.  After multiplying by $L^d$, the quotient bound is
\begin{equation}
 C L^d(1+M)^d(CM/L)^p\|u\|_{r,\bullet},
 \qquad
 (d,p)=(1,2)\ \text{or}\ (1/8,1+1/8).
 \label{rg:marked-quotient}
\end{equation}
Both expressions tend to zero with $L$.  The scaled long-term bound
also tends to zero.  The exact line multiplier is one by
\eqref{rg:line-transport}, proving \eqref{rg:marked-matrix} as before.
Averaging a projection over the finite rotation group preserves its
value on the invariant even line and keeps its norm bounded.

Finally let the input be a boundary family.  A short carrier touching
the edge has its anchor within $CMr$ of it.  \Cref{rad:boundary}, with
support radius $s=CMr$, gives for each such term
\begin{equation}
 \|\pi P_{R,z}f_{A,z}\|_\infty
       \le C(CM/L)^2\|f_{A,z}\|_\infty.
 \label{rg:boundary-term}
\end{equation}
The use of an anchor just off the edge is allowed by that proposition;
one can also recenter at a boundary point at distance at most $s$,
changing only fixed support and separation factors.

For a fixed output rooted square, the relevant anchors lie in a band of
width $CMr$ along an edge segment of length $O(R)$.  In a one-sheet
graph chart this band is covered by at most $CL(1+M)^C$ input
$r$-blocks.  Each block bounds the sum of the input norms anchored
there, since an anchor belongs to its carrier.  Summing
\eqref{rg:boundary-term} therefore costs at most
\[
 CL(1+M)^C(CM/L)^2\|h\|_r=o_L(1)\|h\|_r.
\]
The long terms are bounded by \eqref{rg:long-tail}.  An output square of
a short boundary term still touches the boundary, because the anchor
is within $CMr<R$ of it; a long unchanged carrier retains its contact.
This proves the boundary assertion.
\end{proof}

The three contractions can be arranged simultaneously.  The lower
threshold on $r_0$ depends on the chosen $L$ and on the critical buffer
estimates.  In a fixed bounded domain, the upper permissible physical
radius additionally depends on its boundary: a compact $C^2$ Jordan
curve has uniform one-sheet graph charts and arbitrarily small slope
deviation on sufficiently small charts.  Thus every fixed domain has
the windows needed for all sufficiently fine meshes, without a
domain-dependent choice of $L$, $r_0$, or $\rho$.

The boundary contraction in \Cref{rg:linear} applies to the boundary
family itself.  Leakage from restricting a bulk family is instead
estimated by the bounded map of \Cref{rg:map}; it is not assigned a
spurious contraction factor.  This distinction is used in the next
section when the critical plane trajectory is compared with finite
domains.
